\section{Discussion}

\subsection{Controlled Learning-rate Experiment}\label{controlled_discussion}

From the results obtained from Section \ref{controlled_learningrate_section}, it is clear that changing the learning rate has a clear impact on model training as well as model
prediction. The 1e$^{-4}$ learning rate represents a learning rate that is too small. This is clearly evident from the slowest training time since the model takes very small steps. 
This results in slower convergence and leads to underfitting under the maximum default iterations of a 1000 which can be seen in Table \ref{sgd_learningrate}. On the other 
hand, the 1e$^{-2}$ learning rate represents a learning rate that is too large. The model makes too big steps that may result in the model failing to converge to a minimum. 
As a result, the model oscillates and is unstable. This leads to poor generalisation which can be seen by the worst Test PR-AUC from Table \ref{sgd_learningrate}. Lastly, 
we have the 1e$^{-3}$ learning rate which represents a balanced learning rate. This learning rate produces a model with the best overall results in terms of convergence and 
generalisation as well as training time. This learning rate would be best fit for deployment.

\subsection{Grid Search Hyperparameter Tuning}

From the results obtained from Section \ref{gridsearch_section}, it is evident that there is a clear correlation between these results and the controlled learning rate sweep that
was done in Section \ref{controlled_learningrate_section}. The best grid search configuration had a learning rate of 3e$^{-3}$ which closely aligns with the best learning rate
obtained from the controlled learning rate sweep, 1e$^{-3}$. The best rank configuration also obtained a slightly better validation and test PR-AUC when compared to the best 
learning rate's results obtained from the controlled learning rate sweep. This indicates that other parameters such as alpha and penalty had a role to play since these differed
from the defaults used during the controlled learning rate sweep experiment.

\subsection{Bayesian Optimization Hyperparameter Tuning}\label{bayesian_discussion}

From the results obtained from Section \ref{bayesian_section}, in specific Table \ref{bayesian_grid}, it can be seen that Bayesian optimisation peformed slightly better than
grid search in terms of generalisation. The main reason as to why this is the case boils down to essentially how each method works. Grid Search evaluates every possible 
combination of parameters in the predefined grid. Bayesian optimization on the other hand, creates a probalisitic model of the objective function and uses previous trials to 
decide which parameters to test next. Within the search space, it is able to more intelligently search based on promising previous trials. This is evident in the best 
hyperparameter configuration for both tuning methods shown in Section \ref{gridsearch_section} and Section \ref{bayesian_section}. Both have fairly similar parameters
with the exception that the parameters obtained from Bayesian optimization are more specific which in turn results in a better performance. This is due to the fact that
grid search can only linearly sweep through the grid to find the best combination whereas with Bayesian optimization, the model is able to logarithmically search through
the space with its only constraints being the upper and lower limit of the search space. This generally leads to better overall performance at the cost of computation 
time depending on the number of calls made.

\subsection{Comparative Analysis}

Looking at all the results obtained from the Section \ref{results}, it is clear that there are pros and cons with both grid search and Bayesian optimization hyperparameter tuning methods.
It is clear that when limited compute is of main concern, grid search is preferred over Bayesian optimization. This largely comes down to what was discussed in Section \ref{bayesian_discussion}
since grid search simply evaluates all the combinations whereas Bayesian optimization uses a probalisitic model which is more computationally expensive. However, it must be noted
that this preference largely comes to the search space as well. If a large grid space is required, this can significantly impact computation time and in that scenario, Bayesian
optimization should be favoured instead. If model generalisation is of priority, then Bayesian optimization is preferred. Since Bayesian optimization can look much deeper into
a search space, it can generally obtain a much better hyperparameter configuration that would produce the best evaluation results. This is also evident from the convergence
plot (Figure \ref{convergence_plot}) which indicates that if more calls were used, this optimization technique would be able to find even better configurations that result
in better performance. Lastly, it is also important to consider the learning rate used when optimising. It is clear from Section \ref{controlled_discussion} that the learning
rate play a key role in model stability. Refining the optimum learning rate window at the beginning before running other optimization techniques can significantly reduce
computation time and resources whilst producing better evaluation results.